\clearpage\corrigesChapitre{16}{Lumière et spectres}
\corrigesSouples
\begin{blocLivre}{cor-16-1}\begin{corrige}{1}{Une valeur à connaître}
\begin{questions}
\item $c\simeq\qty{3,00e8}{m.s^{-1}}$.
\item Elle est près de neuf cent mille fois plus grande que la vitesse du son dans l’air.
\end{questions}
\end{corrige}\end{blocLivre}
\begin{blocLivre}{cor-16-2}\begin{corrige}{2}{Des nanomètres aux mètres}
\begin{questions}
\item $\lambda=\qty{5,50e-7}{m}$.
\item Non, c’est une unité de longueur, égale à $10^{-9}$ m.
\end{questions}
\end{corrige}\end{blocLivre}
\begin{blocLivre}{cor-16-3}\begin{corrige}{3}{Continuité du spectre}
\schema[corrige]{spectres}
\begin{questions}
\item La bande couvrant sans interruption un intervalle est continue ; les traits séparés constituent un spectre de raies.
\item Le spectre continu.
\end{questions}
\end{corrige}\end{blocLivre}
\begin{blocLivre}{cor-16-4}\begin{corrige}{4}{Un trajet lumineux}
\begin{questions}
\item $\Delta t=900/(3,00\times10^8)=\qty{3,00e-6}{s}$.
\item $\qty{3,00}{\micro s}$.
\end{questions}
\end{corrige}\end{blocLivre}
\begin{blocLivre}{cor-16-5}\begin{corrige}{5}{Reconnaître des raies}
\begin{questions}
\item A présente le même ensemble de raies à la précision annoncée.
\item Une coïncidence isolée peut être fortuite ou commune à plusieurs espèces ; l’ensemble renforce l’identification.
\end{questions}
\end{corrige}\end{blocLivre}
\begin{blocLivre}{cor-16-6}\begin{corrige}{6}{Chauffer davantage}
\begin{questions}
\item Il s’enrichit vers les radiations de plus courtes longueurs d’onde du visible.
\item Non : un filament dense chaud présente un spectre continu.
\end{questions}
\end{corrige}\end{blocLivre}
\begin{blocLivre}{cor-16-7}\begin{corrige}{7}{Monochromatique ou non}
\begin{questions}
\item Non, elle émet deux radiations distinctes.
\item Non : plusieurs compositions peuvent produire une couleur apparente voisine.
\end{questions}
\end{corrige}\end{blocLivre}
\begin{blocLivre}{cor-16-8}\begin{corrige}{8}{Voir un objet}
\begin{questions}
\item De la lumière de la source est diffusée par le livre puis pénètre dans l’œil.
\item Non, la vision résulte de lumière reçue par l’œil.
\end{questions}
\end{corrige}\end{blocLivre}
\begin{blocLivre}{cor-16-9}\begin{corrige}{9}{Une distance astronomique}
\begin{questions}
\item $d=c\Delta t=\qty{1,50e11}{m}$.
\item $\qty{1,50e8}{km}$. L’observation nous renseigne sur l’état de la source 500 s plus tôt.
\end{questions}
\end{corrige}\end{blocLivre}
\begin{blocLivre}{cor-16-10}\begin{corrige}{10}{Un aller-retour lumineux}
\begin{questions}
\item $d=c\Delta t/2=3,00\times10^8\times40,0\times10^{-9}/2=\qty{6,00}{m}$.
\item La durée mesurée couvre l’aller jusqu’à la cible puis le retour au récepteur.
\end{questions}
\end{corrige}\end{blocLivre}
\begin{blocLivre}{cor-16-11}\begin{corrige}{11}{La lecture d’un mélange}
\begin{questions}
\item La présence de A est compatible avec les raies observées.
\item Non : les raies supplémentaires suggèrent une autre espèce ou demandent une explication expérimentale.
\end{questions}
\end{corrige}\end{blocLivre}
\begin{blocLivre}{cor-16-12}\begin{corrige}{12}{Éclair et tonnerre}
\demarche{Estimer d’abord la distance par le son, puis le temps de parcours lumineux.}
\begin{questions}
\item $d\simeq340\times3,0=\qty{1020}{m}$, soit environ 1,0 km.
\item La lumière parcourt cette distance en environ 3,4 microsecondes, négligeables devant 3,0 s.
\end{questions}
\retiens{Une approximation se justifie par une comparaison d’ordres de grandeur.}
\end{corrige}\end{blocLivre}

